%% Anexo C — Hojas de vida del equipo
%% Fragmento: no se compila solo. Lo incluye docs/entrega/entrega-fase1.tex.

\section{Anexo C. Hojas de vida del equipo}

Se presentan a continuación las hojas de vida de los tres integrantes de
Legasoft, orientadas al rol que cada uno asume en la consultora. Cada una declara
las competencias con que cuenta para cumplir sus funciones, la experiencia y las
certificaciones que las respaldan, y la disponibilidad horaria y el coste
imputable al \lgSemestre, conforme al principio de gestión de recursos humanos
adoptado en la cláusula quinta del acta de constitución.

\begin{nota}[Sobre la valorización del trabajo]
La tarifa interna de Bs 50 por hora es una referencia de la consultora para
estimar el coste del esfuerzo y dimensionar propuestas. No representa
remuneración entre los socios. El coste del semestre resulta de multiplicar la
dedicación semanal por las \lgSemanas\ semanas efectivas del periodo
(\lgPeriodo).
\end{nota}

% ===========================================================================
\subsection{\lgDirector}

\begin{nota}[\lgDirectorRol]
La Paz, Bolivia · \href{mailto:roziel147@gmail.com}{roziel147@gmail.com} ·
\href{https://github.com/ozioziel}{github.com/ozioziel} ·
\href{https://www.linkedin.com/in/oziel-rodman-ramos-torrez-98b914403}{LinkedIn}
\end{nota}

\subsubsection{Perfil profesional}

Estudiante de Ingeniería en Sistemas con perfil orientado a la dirección de
proyectos y el análisis de sistemas. Cuenta con experiencia práctica levantando
necesidades mediante entrevistas, documentando historias de usuario y requisitos
funcionales y no funcionales, y analizando procesos manuales para convertirlos en
propuestas de software. Ha participado en proyectos con clientes y
organizaciones, coordinando la comunicación, la documentación, la validación de
prototipos, la preparación de demostraciones y la entrega de un producto mínimo
viable. Su base técnica en desarrollo full stack, bases de datos, APIs y cloud le
permite comunicarse tanto con usuarios como con equipos de desarrollo.

\subsubsection{Competencias para el rol}

\begin{center}
\begin{tabularx}{\linewidth}{@{}L{4.0cm} Y@{}}
\toprule
\sffamily\bfseries\color{lgPrimario}Ámbito &
\sffamily\bfseries\color{lgPrimario}Competencias \\
\midrule
Análisis de requisitos & Entrevistas; levantamiento de necesidades; requisitos funcionales y no funcionales; especificación de requisitos (SRS); historias de usuario; criterios de aceptación; priorización. \\
Estándares & ISO/IEC/IEEE 29148 como marco de referencia para la especificación, validación y trazabilidad de requisitos, adoptado por la consultora para este rol. \\
Gestión de proyectos & Jira; Scrum; \emph{backlog}; \emph{sprints}; cronograma; asignación y seguimiento de tareas; gestión de entregables. \\
Documentación & SRS; minutas de reunión; informes; definición de alcance; documentación funcional; presentaciones y demostraciones; Swagger/OpenAPI. \\
Modelado de sistemas & BPMN; UML; casos de uso; diagramas de clases y de secuencia; análisis de procesos; arquitectura por capas. \\
Comunicación y liderazgo & Comunicación con clientes; escucha activa; facilitación de reuniones; presentación de propuestas; coordinación de equipos; resolución de problemas. \\
Base técnica & React; TypeScript; Node.js; NestJS; Flutter; PostgreSQL/Supabase; APIs REST; Git y GitHub; Docker; Linux; fundamentos de Cloud. \\
\bottomrule
\end{tabularx}
\captionof{table}{Competencias del \lgDirectorRol.}
\end{center}

\subsubsection{Experiencia relevante}

\begin{description}
  \item[Consultor independiente de software — Ipas Bolivia (abril 2026)]
  Consultoría remunerada para formalizar, ajustar y entregar un producto mínimo
  viable surgido de una hackatón. Apoyó la coordinación con la contraparte, la
  organización técnica del trabajo y el seguimiento de los entregables acordados;
  documentó el funcionamiento del producto y colaboró en la demostración, la
  presentación funcional y el informe de entrega. El código y los materiales
  internos permanecen privados por obligaciones de confidencialidad.

  \item[Desarrollador en proyecto colaborativo — NexoLab (junio 2026 a la fecha)]
  Integrante del equipo ganador del reto Hipermaxi en la InnovaHack Santa Cruz
  2026. Analizó las restricciones reales de integración y las necesidades del
  entorno empresarial para adaptar el prototipo inicial, y colaboró en su
  evolución hacia un \emph{widget} embebible con arquitectura desacoplada,
  facilitando la coordinación entre interfaz, backend y componentes de
  inteligencia artificial.
\end{description}

\subsubsection{Proyectos de análisis y levantamiento de requisitos}

\begin{itemize}
  \item \textbf{MAPIA — análisis y diseño de aplicación geolocalizada (2026).}
        Definición de módulos, flujos de usuario y reglas para una aplicación
        móvil con autenticación, mapas, geolocalización, almacenamiento de
        imágenes e integración de servicios. La documentación de la API mediante
        Swagger y la separación frontend-backend facilitaron la coordinación y la
        trazabilidad técnica. Flutter, NestJS, PostgreSQL/PostGIS, Google Maps.
  \item \textbf{Prototipo de acceso a la justicia y evidencia digital
        (julio–agosto 2026).} En una hackatón de la Fundación CONSTRUIR,
        participó en el análisis de un problema social complejo y en la
        definición de un flujo guiado para registrar evidencia, consultar rutas
        de denuncia y localizar centros de apoyo. El proyecto exigió organizar
        requisitos funcionales, información legal y criterios de privacidad en
        una propuesta comprensible para usuarios y equipo técnico.
\end{itemize}

\subsubsection{Logros y trabajo en equipo}

\begin{itemize}
  \item Primer lugar en la InnovaHack Santa Cruz 2026, con el equipo NexoLab en
        el reto Hipermaxi (junio 2026).
  \item Segundo lugar en la Hackatón de Innovación Digital de Ipas Bolivia (marzo
        2026); la propuesta fue seleccionada para una consultoría profesional
        posterior.
  \item Segundo lugar en la Feria FUNCTEC de la Universidad Católica Boliviana
        (septiembre 2024).
\end{itemize}

\subsubsection{Formación y certificaciones}

Ingeniería en Sistemas, Universidad Católica Boliviana «San Pablo», sede La Paz
(febrero 2024 a la fecha): análisis y diseño de sistemas, ingeniería de
requisitos, modelado BPMN y UML, bases de datos, algoritmos, tecnologías web,
desarrollo móvil, arquitectura de software y fundamentos de gestión tecnológica.

\begin{itemize}
  \item \textbf{Google Cloud Data Analytics Certificate} — análisis de datos y
        servicios de Google Cloud; credencial verificable en Credly.
  \item \textbf{Microsoft 365 Copilot Skilling Campaign 2026} — uso de IA
        generativa para documentación y productividad.
  \item \textbf{Networking Basics} — fundamentos de redes y comunicación entre
        sistemas.
\end{itemize}

Idiomas: español nativo; inglés \emph{upper-intermediate}, con solvencia
conversacional y técnica para lectura de documentación.

\subsubsection{Disponibilidad y coste}

\begin{center}
\begin{tabularx}{\linewidth}{@{}L{4.6cm} Y@{}}
\toprule
\sffamily\bfseries\color{lgPrimario}Concepto &
\sffamily\bfseries\color{lgPrimario}Compromiso para el \lgSemestre \\
\midrule
Dedicación semanal & 12 horas efectivas. \\
Franjas horarias & Lunes, miércoles y viernes de 19:00 a 22:00; sábados de 09:00 a 12:00. \\
Modalidad & Trabajo remoto asincrónico sobre el repositorio y Jira; reuniones presenciales o virtuales según convenga al cliente. \\
Horas del semestre & 192 horas (\lgSemanas\ semanas efectivas). \\
Tarifa interna & Bs 50 por hora. \\
Coste imputable & Bs 9.600. \\
\bottomrule
\end{tabularx}
\captionof{table}{Disponibilidad y coste del \lgDirectorRol.}
\end{center}

% ===========================================================================
\clearpage
\subsection{\lgArquitecto}

\begin{nota}[\lgArquitectoRol]
La Paz, Bolivia ·
\href{mailto:luisdanielrojascaceres@gmail.com}{luisdanielrojascaceres@gmail.com} ·
\href{https://github.com/daniel69zz}{github.com/daniel69zz}
\end{nota}

\subsubsection{Perfil profesional}

Estudiante de Ingeniería en Sistemas en la Universidad Católica Boliviana «San
Pablo», con interés en el desarrollo de software, las bases de datos, el backend,
el frontend, la computación en la nube y la inteligencia artificial. Cuenta con
experiencia académica y práctica desarrollando aplicaciones web, sistemas
backend, videojuegos y soluciones tecnológicas. Posee conocimientos sólidos en
diseño, modelado y administración de bases de datos relacionales, incluyendo SQL
avanzado, disparadores, consultas complejas, relaciones, restricciones e
integridad de datos. Está certificado como AWS Certified Cloud Practitioner.

\subsubsection{Competencias para el rol}

\begin{center}
\begin{tabularx}{\linewidth}{@{}L{4.0cm} Y@{}}
\toprule
\sffamily\bfseries\color{lgPrimario}Ámbito &
\sffamily\bfseries\color{lgPrimario}Competencias \\
\midrule
Arquitectura & Modelo C4 para la representación de contexto, contenedores y componentes; patrones de diseño; principios SOLID; arquitectura cliente-servidor y basada en microservicios; documentación de decisiones (ADR). \\
Backend & Node.js; NestJS; Express; Spring Boot; diseño e implementación de APIs REST; JWT; Swagger; TypeORM. \\
Bases de datos & PostgreSQL; MySQL; \textbf{SQL Server}; Supabase; modelado relacional y normalización. \\
SQL avanzado & Consultas complejas; \emph{joins}; subconsultas; disparadores; funciones; secuencias; vistas; restricciones; llaves primarias y foráneas; auditoría e integridad referencial. \\
Lenguajes & JavaScript; TypeScript; Python; C++; SQL. \\
Frontend & React; Vite; HTML; CSS; Tailwind CSS. \\
DevOps y cloud & Docker; Git y GitHub; Linux; AWS; Oracle Cloud. \\
\bottomrule
\end{tabularx}
\captionof{table}{Competencias del \lgArquitectoRol.}
\end{center}

\subsubsection{Proyectos destacados}

\begin{description}
  \item[Solución de IA para soporte a proveedores — Hipermaxi (2026)]
  Participación con el equipo NexoLab en la InnovaHack 2026 de CAINCO, obteniendo
  el primer lugar general y el primer lugar en el Reto Empresarial. Desarrollo de
  una propuesta basada en inteligencia artificial para apoyar el servicio,
  soporte y guía de proveedores dentro de la plataforma de Hipermaxi, aplicando
  conceptos de desarrollo de software, APIs, automatización y experiencia de
  usuario en un contexto empresarial real.

  \item[Sistema de Gestión Deportiva Universitaria (2026)]
  Sistema web para la gestión de actividades deportivas universitarias: usuarios,
  roles, equipos, torneos, reservas y reportes. Diseño y creación de la base de
  datos para el Departamento de Deportes de la Universidad Católica Boliviana
  «San Pablo», con un modelo relacional unificado para el consumo de todos los
  microservicios y aplicación de SQL avanzado —consultas complejas,
  normalización, integridad referencial, disparadores y secuencias—. NestJS,
  React, PostgreSQL, Docker, TypeScript y APIs REST.

  \item[Aplicación de evidencias y soporte legal (2025)]
  Desarrollo de una aplicación para la gestión de evidencias y asistencia legal,
  integrando tecnologías web y herramientas de inteligencia artificial para
  apoyar el análisis de incidentes. Segundo lugar en el Hackathon IPAS Bolivia.

  \item[Sistema mini e-commerce y sistema TPS para cafetería (2025)]
  Backend en Node.js para la gestión de productos, usuarios y operaciones, con
  APIs REST, estructura de rutas, controladores y manejo de datos; y un sistema
  de procesamiento de transacciones con frontend y backend para el registro de
  pedidos y el control de operaciones del negocio.

  \item[Modelo de reconocimiento de lenguaje de señas y videojuego en C++ (2025)]
  Proyecto de aprendizaje automático para el reconocimiento de gestos con
  TensorFlow, OpenCV y MediaPipe; y un videojuego inspirado en Pac-Man
  desarrollado desde cero en C++, aplicando estructuras de datos y programación
  orientada a objetos.
\end{description}

\subsubsection{Formación, certificaciones y logros}

Ingeniería en Sistemas, Universidad Católica Boliviana «San Pablo», La Paz (2024
a la fecha). Informática, Universidad Mayor de San Andrés, La Paz (2024 a la
fecha).

\begin{itemize}
  \item \textbf{AWS Certified Cloud Practitioner} — Amazon Web Services.
  \item Primer lugar general y primer lugar en el Reto Empresarial de la
        InnovaHack 2026 de CAINCO, con el equipo NexoLab.
  \item Segundo lugar en el Hackathon IPAS Bolivia.
  \item Segundo lugar en la feria tecnológica universitaria FUNTEC-UCB.
  \item Diseño y desarrollo de la base de datos del sistema del Departamento de
        Deportes de la Universidad Católica Boliviana «San Pablo».
\end{itemize}

Idiomas: español nativo; inglés intermedio.

\subsubsection{Disponibilidad y coste}

\begin{center}
\begin{tabularx}{\linewidth}{@{}L{4.6cm} Y@{}}
\toprule
\sffamily\bfseries\color{lgPrimario}Concepto &
\sffamily\bfseries\color{lgPrimario}Compromiso para el \lgSemestre \\
\midrule
Dedicación semanal & 12 horas efectivas. \\
Franjas horarias & Martes, jueves y viernes de 19:00 a 22:00; sábados de 09:00 a 12:00. \\
Modalidad & Trabajo remoto asincrónico sobre el repositorio y Jira; disponibilidad para revisiones técnicas y sesiones de integración. \\
Horas del semestre & 192 horas (\lgSemanas\ semanas efectivas). \\
Tarifa interna & Bs 50 por hora. \\
Coste imputable & Bs 9.600. \\
\bottomrule
\end{tabularx}
\captionof{table}{Disponibilidad y coste del \lgArquitectoRol.}
\end{center}

% ===========================================================================
\clearpage
\subsection{\lgQA}

\begin{nota}[\lgQARol]
La Paz, Bolivia ·
\href{mailto:pedroandrez.paca@gmail.com}{pedroandrez.paca@gmail.com}
\end{nota}

\subsubsection{Perfil profesional}

Estudiante de quinto semestre de Ingeniería de Sistemas, con promedio académico
superior al 80\%. Se distingue por su responsabilidad, orden y disciplina en el
cumplimiento de sus compromisos académicos, cualidades directamente aplicables al
aseguramiento de la calidad, donde el valor del rol depende de la constancia en
la ejecución de las pruebas y en el registro de la evidencia. Posee formación en
fundamentos de programación, bases de datos y sistemas, y experiencia sostenida
de apoyo técnico a sus pares.

\subsubsection{Competencias actuales}

\begin{center}
\begin{tabularx}{\linewidth}{@{}L{4.0cm} Y@{}}
\toprule
\sffamily\bfseries\color{lgPrimario}Ámbito &
\sffamily\bfseries\color{lgPrimario}Competencias \\
\midrule
Programación y datos & Fundamentos de programación y de bases de datos. \\
Entornos de desarrollo & Manejo de entornos de desarrollo y herramientas de programación. \\
Sistemas operativos & Instalación y configuración de sistemas operativos; preparación de entornos de trabajo. \\
Equipos y redes & Ensamblaje y configuración inicial de equipos; configuración básica de redes locales y parámetros de conexión. \\
Ofimática y documentación & Manejo de herramientas ofimáticas para la elaboración y el registro de documentación. \\
Cualidades para el rol & Orden, disciplina y constancia en el seguimiento de procedimientos; atención al detalle; disposición para el trabajo en equipo. \\
\bottomrule
\end{tabularx}
\captionof{table}{Competencias actuales del \lgQARol.}
\end{center}

\subsubsection{Plan de habilitación para el rol}

La consultora aplica el valor de transparencia también a la evaluación de sus
propias capacidades. El perfil actual acredita fundamentos sólidos, pero no
todavía la práctica específica que exigen el aseguramiento de calidad y el
desarrollo frontend. Por ello el equipo acuerda el siguiente plan de
habilitación, cuyo cumplimiento se verifica en las revisiones internas del
semestre:

\begin{center}
\begin{tabularx}{\linewidth}{@{}L{4.4cm} Y C{2.2cm}@{}}
\toprule
\sffamily\bfseries\color{lgPrimario}Competencia a incorporar &
\sffamily\bfseries\color{lgPrimario}Evidencia esperada &
\sffamily\bfseries\color{lgPrimario}Plazo \\
\midrule
Diseño de casos de prueba y gestión de defectos & Plan de pruebas y casos derivados de los criterios de aceptación del SRS, en \ruta{docs/calidad/}. & Antes del hito H1 \\
ISO/IEC 25010 & Selección justificada de los atributos de calidad a evaluar en el producto. & Antes del hito H2 \\
Desarrollo frontend (web o Flutter) & Componentes de interfaz versionados en \ruta{src/frontend/}. & Antes del hito H3 \\
Automatización y CI/CD & Flujo de GitHub Actions que ejecute las pruebas en cada integración. & Antes del hito H3 \\
\bottomrule
\end{tabularx}
\captionof{table}{Plan de habilitación acordado para el \lgQARol.}
\end{center}

El acompañamiento técnico corresponde al \lgArquitectoRol, conforme a la relación
entre roles descrita en \texttt{LGS-INST-001}.

\subsubsection{Experiencia}

\begin{description}
  \item[Apoyo técnico académico — independiente (2021 a la fecha)]
  Asistencia a compañeros en la resolución de ejercicios y problemas de
  programación y sistemas. Instalación y configuración de sistemas operativos y
  software académico; ensamblaje y configuración inicial de equipos informáticos;
  configuración básica de redes locales; uso constante de entornos de desarrollo
  y herramientas informáticas en actividades académicas.
\end{description}

\subsubsection{Formación y actividades}

Ingeniería de Sistemas, Universidad Católica Boliviana «San Pablo», La Paz
(febrero 2024 a la fecha), estudiante regular de quinto semestre. Técnico Medio
en Sistemas Informáticos, Unidad Educativa AMERINST, La Paz (febrero 2022 a
diciembre 2023).

Participación activa en el Ballet Folclórico de la Universidad Católica Boliviana
«San Pablo», actividad en la que ha desarrollado disciplina, compromiso y trabajo
en equipo.

Idiomas: español nativo; inglés de nivel intermedio, con capacidad de lectura
técnica.

\subsubsection{Disponibilidad y coste}

\begin{center}
\begin{tabularx}{\linewidth}{@{}L{4.6cm} Y@{}}
\toprule
\sffamily\bfseries\color{lgPrimario}Concepto &
\sffamily\bfseries\color{lgPrimario}Compromiso para el \lgSemestre \\
\midrule
Dedicación semanal & 12 horas efectivas, de las cuales 2 se destinan al plan de habilitación. \\
Franjas horarias & Lunes, martes y jueves de 19:00 a 22:00; sábados de 09:00 a 12:00. \\
Modalidad & Trabajo remoto asincrónico sobre el repositorio y Jira; ejecución de pruebas y registro de evidencia antes de cada entrega. \\
Horas del semestre & 192 horas (\lgSemanas\ semanas efectivas). \\
Tarifa interna & Bs 50 por hora. \\
Coste imputable & Bs 9.600. \\
\bottomrule
\end{tabularx}
\captionof{table}{Disponibilidad y coste del \lgQARol.}
\end{center}

% ===========================================================================
\subsection{Resumen de recursos humanos}

La franja del sábado de 09:00 a 12:00 es común a los tres integrantes y se
reserva para la reunión semanal de coordinación y para las reuniones con la
organización cliente.

\begin{center}
\begin{tabularx}{\linewidth}{@{}Y C{2.2cm} C{2.2cm} C{2.4cm}@{}}
\toprule
\sffamily\bfseries\color{lgPrimario}Integrante y rol &
\sffamily\bfseries\color{lgPrimario}Horas por semana &
\sffamily\bfseries\color{lgPrimario}Horas del semestre &
\sffamily\bfseries\color{lgPrimario}Coste imputable \\
\midrule
\lgDirector \newline \textit{\small \lgDirectorRol}     & 12 & 192 & Bs 9.600 \\
\lgArquitecto \newline \textit{\small \lgArquitectoRol} & 12 & 192 & Bs 9.600 \\
\lgQA \newline \textit{\small \lgQARol}                 & 12 & 192 & Bs 9.600 \\
\midrule
\textbf{Capacidad total de la consultora} & \textbf{36} & \textbf{576} & \textbf{Bs 28.800} \\
\bottomrule
\end{tabularx}
\captionof{table}{Capacidad y coste del equipo para el \lgSemestre.}
\end{center}

Esta capacidad —36 horas semanales— es la restricción de recursos con la que se
dimensionará el alcance comprometido con la organización cliente. Cualquier
solicitud que la exceda deberá traducirse en una reducción de alcance o en una
extensión de plazo, y no en una sobrecarga no declarada del equipo.
